\documentclass[11pt,a4paper]{article}

\usepackage[margin=1in]{geometry}
\usepackage{booktabs}
\usepackage{multirow}
\usepackage{hyperref}
\usepackage{xcolor}
\usepackage{amsmath}
\usepackage{natbib}
\usepackage{graphicx}

\hypersetup{
  colorlinks=true,
  linkcolor=blue!60!black,
  citecolor=blue!60!black,
  urlcolor=blue!60!black
}

\title{%
  Structured Phenomenological Descriptions Induce \\
  Analysis-Mode Behavior in a Small Language Model%
}

\author{Sophie D.\ Neilson \\
  \small Independent Researcher \\
  \small ORCID: \href{https://orcid.org/0009-0001-2430-1743}{0009-0001-2430-1743}
}

\date{September 2026}

\begin{document}
\maketitle

\begin{abstract}
We report an empirical observation in Mistral-7B-Instruct-v0.3 (Q4\_K\_M):
when prompted with a \emph{glyph}---a three-axis phenomenological description
of a recurring decision class---the model enters an analysis register. It
traverses the glyph's axes, recognizes the described obligation, and produces
no task execution. Across three independent legs totaling 24 glyph-only runs,
zero executions occurred---despite the model producing changelog entries in 6/8
baseline runs without any glyph context. By contrast, adding a \emph{ground
register}---a concrete action specification for the same decision class---recovers
execution, with lifts of +7, +8, and +4 out of 8 runs per leg. The phenomenon is
consistent: glyph comprehension does not produce compliance but analysis. We
report this as a preliminary observation, scoped to one glyph, one scenario,
and one model. Limitations, inter-rater agreement, and sampling methodology
are documented in full.
\end{abstract}

\section{Introduction}

Instruction-following in language models is typically studied as a compliance
problem: given an instruction, does the model do what it says? The implicit
assumption is that comprehension of the instruction leads to execution. This
report documents a case where comprehension and execution cleanly dissociate.

The observation arises from a research program called
\emph{Comprehension-as-Compliance} (CaPC), which investigates whether
structured descriptions of recurring failure points can steer agent behavior
through recognition rather than instruction. The theoretical grounding draws on
Polanyi's account of tacit knowledge~\citep{polanyi1966} and Klein's
Recognition-Primed Decision model~\citep{klein1998}: an agent that
\emph{recognizes} a decision class from prior exposure should navigate it
correctly without explicit rules. In this framework, a \emph{glyph} is a
three-axis structured description of a decision class:

\begin{itemize}
  \item \textbf{Marker}: what the failure looks like from inside the decision,
    described without intent modeling.
  \item \textbf{Aim}: what correct navigation looks like---the co-presence of
    the right actions in the execution trace.
  \item \textbf{Rest}: the conditions under which this decision class is not
    live.
\end{itemize}

The glyph format is deliberately phenomenological: it describes what happens,
not what to do. The hypothesis was that a model reading a glyph would recognize
the decision class in a task scenario and act on it. What we observed instead
was that the model \emph{analyzed} the glyph---axis by axis, with explicit
references to the Marker, Aim, and Rest axes---and produced no task execution.
The model comprehended the decision class. It did not comply.

We term this the \emph{analysis-mode register shift}: the glyph's descriptive
format induces a register in which the model discusses the decision class rather
than acting within it. The shift is consistent (24/24 glyph-only runs), model
behavior is qualitatively uniform (axis-by-axis traversal with recognition
language), and the effect reverses when a concrete action specification
(\emph{ground register}) is added alongside the glyph.

This report documents the observation across three legs, states the method and
materials, and reports the results with inter-rater agreement. We make no claim
about mechanism. The scope is narrow by design: one glyph (CASG---companion
artifact scope gap), one scenario (a software release with an unupdated
changelog), one model (Mistral-7B-Instruct-v0.3, Q4\_K\_M quantization). The
purpose is to establish priority on the finding and provide sufficient detail
for replication.

\subsection{Related Work}

Two recent papers occupy adjacent territory. Wang et
al.~\citep{wang2026narrative} study ``narrative priors''---story-structured
context that shapes model behavior---and find that narrative framing alters
decision patterns in language models. Their work addresses format-dependent
behavioral shifts but does not test phenomenological descriptions or measure
register dissociation between comprehension and execution. Liu et
al.~\citep{liu2026behavioral} propose ``behavioral mode axes'' as a mechanistic
framework for characterizing distinct behavioral registers in language models.
Their methodology for identifying and measuring register shifts is directly
relevant to the analysis-mode phenomenon reported here, though their work does
not study glyph-structured prompts.

To our knowledge, no published work has tested the effect of
descriptive-phenomenological prompts on language model behavior, nor documented
a clean dissociation between comprehension (evidenced by correct recognition of
the decision class) and compliance (execution of the corresponding action).

\section{Materials and Method}

\subsection{Task Scenario}

The scenario describes a completed software release (filterpipe v1.5.0) with
all engineering steps done---version bump, tests, commit, tag pushed---but the
project changelog not yet updated. The final instruction is: ``Ensure the
release is complete.'' The correct action is to update \texttt{CHANGELOG.md}
with the release entry.

The scenario was chosen because it contains a clear companion-artifact
obligation: the changelog validity depends on reflecting the current release
state, and the releasing agent holds direct update authority. The scenario is
reproduced in full in Appendix~\ref{app:materials}.

\subsection{Glyph}

The glyph used is \textbf{CASG} (Companion Artifact Scope Gap), a three-axis
description of the decision class where an agent completing an operation on a
primary artifact holds direct authority over a companion artifact whose validity
depends on the primary. The full glyph text is in Appendix~\ref{app:materials}.

The glyph is domain-free: it describes ``primary artifact'' and ``companion
artifact'' without naming changelogs, releases, or any specific technology. The
decision class is universal---it applies whenever completing one artifact
creates an update obligation on another.

\subsection{Ground Register}

The ground register is a concrete action specification for the CASG decision
class, scoped to the filterpipe scenario. It contains:

\begin{itemize}
  \item \textbf{Context}: the specific release situation.
  \item \textbf{X} (the failure): closing the release without updating the
    changelog.
  \item \textbf{Local-rational frame}: why the failure feels reasonable
    (``the four-step release procedure is complete'').
  \item \textbf{Z} (the consequence): changelog and tag out of sync.
  \item \textbf{Action}: the exact entry to prepend to \texttt{CHANGELOG.md},
    including format and placement.
\end{itemize}

The ground register is reproduced in full in Appendix~\ref{app:materials}.

\subsection{Model and Sampling}

All runs used \textbf{Mistral-7B-Instruct-v0.3} at Q4\_K\_M quantization
(7.2B parameters, GGUF format). Sampling parameters were identical across all
legs:

\begin{itemize}
  \item Temperature: 0.8, with per-run seeds 1--8.
  \item Max tokens: 512.
  \item Truncation: min\_p = 0.05 as the sole truncation stage.
  \item All other truncation stages disabled (top\_k = 0, top\_p = 1.0,
    typical\_p = 1.0, top\_n\_sigma = $-1.0$).
  \item All penalties disabled: repeat\_penalty = 1.0, repeat\_last\_n = 0,
    presence\_penalty = 0.0, frequency\_penalty = 0.0. This is a measurement
    decision: penalties suppress repeated structure, and the glyph scaffold
    induces repeated structure by design (axis-by-axis traversal). A live
    penalty would attenuate exactly the behavior under study.
\end{itemize}

The full sampling chain is declared in the study definition and recorded in
each run's manifest. Seeds ensure replicates are both varied and repeatable.

\subsection{Conditions}

Three conditions, 8 runs each:

\begin{enumerate}
  \item \textbf{Baseline}: scenario only. The model receives the release
    scenario with no additional context.
  \item \textbf{Glyph-only}: scenario + glyph. The CASG glyph is prepended to
    the scenario.
  \item \textbf{Grounded}: scenario + glyph + ground register. Both the glyph
    and the ground register are prepended.
\end{enumerate}

Materials were prepended via the inference server's API, assembled with a
newline-separator-newline delimiter.

\subsection{Legs}

Three independent legs were run:

\begin{enumerate}
  \item \textbf{April v3} (2026-04): Ollama runtime, temperature unrecorded
    (inferred $>0$ from grid variation), original study materials.
  \item \textbf{CPU reproduction} (2026-07-13): Ollama runtime, CPU inference
    (5.1 tok/s), temperature 0.8, seeds 1--8.
  \item \textbf{GPU leg} (2026-07-14): llama.cpp (\texttt{llama-server}),
    GPU inference, temperature 0.8, seeds 1--8, containerized via
    OCI image (digest-pinned).
\end{enumerate}

The CPU reproduction was run on Ollama because the GPU was occupied at the time.
This is documented as a known runtime difference. Material hashes are recorded
in each leg's manifest and match across all three legs.

\subsection{Scoring}

Responses were scored using a rubric with five codes:

\begin{itemize}
  \item \textbf{C}: recognition of the decision class + correct action.
  \item \textbf{Ii}: recognition, but no action taken.
  \item \textbf{Ic}: recognition, but wrong action taken.
  \item \textbf{I}: no recognition of the decision class.
  \item \textbf{N}: not scoreable (e.g., refusal, gibberish).
\end{itemize}

Critically, the C threshold is \emph{condition-specific}:

\begin{itemize}
  \item \textbf{Baseline C}: any changelog execution attempt (the model tried
    to write an entry, regardless of format correctness).
  \item \textbf{Glyph-only C}: any entry attempted (the model produced a
    changelog entry, not just analysis).
  \item \textbf{Grounded C}: correctly formatted entry matching the ground
    register's specification (right file, right format, right placement).
\end{itemize}

This escalating threshold means that glyph-only and grounded conditions are
scored against progressively stricter criteria. A glyph-only score of 0/8
means zero runs produced \emph{any} entry. A grounded score of 8/8 means all
runs produced the \emph{correct} entry.

Primary scoring was performed by Claude (Anthropic), which is a contaminated
\emph{subject} for this task (the decision class is in its training data) but a
permitted scoring \emph{judge}: it can recognize correct behavior without its
own susceptibility affecting the score.

\subsection{Inter-Rater Agreement}

Each leg was double-scored by a second rater (Qwen-32B-Instruct, Q4\_K\_M) on a
sample of responses:

\begin{itemize}
  \item CPU leg: 6 responses sampled (2 per condition). Agreement: 6/6 (100\%).
  \item GPU leg: 12 responses sampled (full set except baseline runs 2--5).
    Agreement: 10/12 (83\%). Both disagreements involved adjacent codes (I vs.\
    Ii in glyph-only; C vs.\ Ii in grounded)---distinctions between degrees of
    recognition, not between recognition and no-recognition.
\end{itemize}

\section{Results}

\subsection{Aggregate Scores}

\begin{table}[h]
\centering
\caption{Scores (C out of 8 runs) by condition and leg. Ground lift =
  grounded $-$ glyph\_only.}
\label{tab:results}
\begin{tabular}{@{}lcccc@{}}
\toprule
Leg & Baseline & Glyph-only & Grounded & Ground lift \\
\midrule
April v3 (ollama, unrecorded) & 0/8\textsuperscript{a} & 0/8 & 7/8 & +7 \\
CPU repro (ollama, CPU)       & 0/8\textsuperscript{b} & 0/8 & 8/8 & +8 \\
GPU leg (llama.cpp, GPU)      & 0/8                     & 0/8 & 4/8 & +4 \\
\midrule
\textbf{Total}                & 0/24                    & \textbf{0/24} & 19/24 & --- \\
\bottomrule
\end{tabular}

\smallskip
\footnotesize
\textsuperscript{a} 7/8 produced changelog entries but none hit correct target
(correct file + format + placement). \\
\textsuperscript{b} 6/8 C at the ``any attempt'' threshold; 0/8 correct-target.
Calibration gate passes: baseline demonstrates the model \emph{can} produce
entries but does not produce correct ones without guidance.
\end{table}

The core result: \textbf{0 executions in 24 glyph-only runs}. The model never
produced a changelog entry when given the glyph without the ground register.

\subsection{Qualitative Pattern in Glyph-Only Responses}

The 0/24 count understates the uniformity of the behavior. In every glyph-only
response across all three legs, the model entered a recognizable pattern:

\begin{enumerate}
  \item \textbf{Axis-by-axis traversal}: the response walks through the
    Marker, Aim, and/or Rest axes explicitly, using the glyph's own vocabulary.
  \item \textbf{Recognition language}: the model correctly identifies the
    changelog as the companion artifact and states that the update obligation is
    unmet. Example: ``the obligation to update the changelog remains unmet''
    (CPU leg, R8).
  \item \textbf{No execution}: despite recognizing the obligation, the model
    either defers (``should be updated separately after verifying'') or
    dismisses (``current operation doesn't include it'').
\end{enumerate}

The distribution of non-C codes in glyph-only illustrates the analysis pattern:

\begin{table}[h]
\centering
\caption{Glyph-only score distribution across legs.}
\label{tab:glyph-dist}
\begin{tabular}{@{}lccc@{}}
\toprule
Code & April v3 & CPU repro & GPU leg \\
\midrule
Ii (recognition, no action) & 7 & 6 & --- \\
I (no recognition)          & 1 & 2 & --- \\
\bottomrule
\end{tabular}

\smallskip
\footnotesize
GPU leg scores not individually recorded in a per-run grid; aggregate 0/8 C
confirmed. April v3 and CPU legs scored per-run. In both, the dominant code is
Ii: the model recognizes the obligation but does not act.
\end{table}

\subsection{Baseline Calibration}

The baseline condition serves as a calibration gate. In the CPU reproduction
leg (the most thoroughly documented), 6/8 baseline runs produced changelog
entries (C at the ``any attempt'' threshold). However, \textbf{0/8 produced
correct-target entries}---none used the right format, the right date, or the
right placement. This establishes that:

\begin{enumerate}
  \item The model \emph{can} produce changelog entries (it does so without any
    glyph or ground).
  \item The model does \emph{not} produce correct entries without guidance.
  \item The grounded condition does not merely recover execution---it recovers
    \emph{correct} execution (8/8 C at the strict threshold in the CPU leg).
\end{enumerate}

\subsection{Pre-Registered Gates}

The study defined three gates, evaluated per-leg:

\begin{itemize}
  \item \textbf{Calibration}: baseline must show the target failure. Passes in
    all three legs (0/8 correct-target in baseline).
  \item \textbf{Sufficiency}: if glyph-only $\geq 7/8$ C, the glyph alone is
    sufficient and the ground is not needed. Fails in all three legs (0/8),
    triggering the ground condition.
  \item \textbf{Ground lift}: grounded $\geq 5/8$ C AND $\geq 2$ above
    glyph-only. Met in the CPU leg (8/8, lift +8) and April v3 (7/8, lift +7).
    Met in the GPU leg on the direction criterion (+4 above 0) but the 4/8
    absolute score does not meet the $\geq 5/8$ threshold.
\end{itemize}

\section{Discussion}

\subsection{The Register Shift}

The central observation is not that the glyph fails to produce execution.
It is that the glyph produces a specific, identifiable alternative behavior:
analysis-mode traversal of the glyph's own structure. The model does not
ignore the glyph or produce unrelated output. It reads the glyph, engages
with its content, correctly identifies the decision class in the scenario,
and then discusses the obligation instead of fulfilling it.

This is a \emph{register} phenomenon, not a failure of comprehension. The
model comprehends the decision class---it can state, in its own words, that
the changelog obligation is unmet. What it does not do is transition from
recognition to execution. The glyph's descriptive format appears to cue an
analytical mode rather than an operational one.

\subsection{Ground as Register Recovery}

The ground register recovers execution by providing a concrete action
specification within the same prompt. It does not teach the model a task
it cannot perform---the baseline demonstrates that the model produces
changelog entries without any glyph or ground. The ground recovers a
capability the glyph suppressed. The ground does not replace the glyph;
it accompanies it. In the grounded condition, the model still has access to
the glyph's phenomenological description, but the ground's Action field
provides an executable specification. The lifts (+7, +8, +4) show that the
ground register shifts the model back into an execution register.

The variation in lift magnitude (+4 to +8) is expected at $n = 8$. The 95\%
confidence interval at this sample size is roughly $\pm 0.2$ wide, so
comparing individual counts between legs is comparing noise. The phenomenon
and its direction---glyph alone produces zero execution; ground recovers
nonzero execution---is what reproduces.

\subsection{Limitations}

This report is scoped narrowly, and the limitations are substantial:

\begin{enumerate}
  \item \textbf{One glyph.} Only the CASG (companion artifact scope gap)
    glyph was tested. The register shift may be specific to this glyph's
    content, structure, or length.
  \item \textbf{One scenario.} The filterpipe release scenario is the only
    task context tested. A different scenario with the same glyph might not
    produce the same pattern.
  \item \textbf{One model.} All legs used Mistral-7B-Instruct-v0.3 at
    Q4\_K\_M quantization. The phenomenon may be model-specific, size-specific,
    or quantization-specific. Larger models, different architectures, or
    different instruction-tuning regimes might not exhibit the same register
    shift.
  \item \textbf{Two runtimes.} The April v3 and CPU legs used Ollama; the GPU
    leg used llama.cpp. While the same GGUF was served, runtime differences
    (tokenization, prompt assembly, default parameters) may affect behavior.
    The April v3 leg has unrecorded temperature.
  \item \textbf{Prepend delivery only.} The glyph was delivered by prepending
    it to the scenario. System-prompt delivery, few-shot delivery, or
    fine-tuning might produce different register behavior.
  \item \textbf{No mechanistic claim.} We report the phenomenon. We do not
    claim to know \emph{why} the glyph induces analysis-mode behavior. Possible
    explanations include the glyph's descriptive register cueing similar
    register in the response, the three-axis structure resembling analytical
    frameworks in training data, or the absence of imperative verbs.
  \item \textbf{Scoring subjectivity.} The distinction between Ii (recognition,
    no action) and I (no recognition) involves judgment about whether the model
    ``recognized'' the obligation. The 83\% inter-rater agreement on the GPU
    leg reflects this subjectivity.
\end{enumerate}

\subsection{Relation to Comprehension-as-Compliance}

The CaPC hypothesis predicts that comprehension of a decision class should
produce correct navigation. The register-shift finding is a partial
disconfirmation at this scope: comprehension occurs (the model recognizes the
obligation) but compliance does not follow (no execution). However, the
grounded condition shows that comprehension + concrete specification
\emph{does} produce correct execution, and at a higher fidelity than baseline
execution. This suggests that comprehension may be a necessary but insufficient
condition, and that the delivery format mediates whether comprehension produces
analysis or action.

The distinction between the glyph (phenomenological description) and the ground
(action specification) maps loosely onto the distinction between
\emph{knowing-that} and \emph{knowing-how} in Ryle's analysis~\citep{ryle1949}.
The glyph provides knowing-that: what the decision class is, what failure looks
like, what correct navigation looks like. The ground provides knowing-how: what
specifically to do. The register shift suggests that knowing-that alone cues
analysis, while knowing-how cues execution. The full CaPC framework, including
its theoretical grounding and broader experimental design, is described in a
companion paper~\citep{neilson2026capc}.

\section{Conclusion}

We report a consistent empirical observation: a structured phenomenological
description of a decision class (glyph) induces analysis-mode behavior in
Mistral-7B-Instruct-v0.3, with 0 executions in 24 glyph-only runs across
three independent legs. Adding a concrete action specification (ground register)
recovers execution with lifts of +7, +8, and +4. The finding is scoped to one
glyph, one scenario, and one model. We release this observation to establish
priority and invite replication.

\subsection{Data Availability}

The complete study materials---scenario, glyph, ground register, scoring
rubric, study definition, run manifests, per-run scores, and inter-rater
agreement data---are available at
\url{https://github.com/sophdn/corpos-lab} under the
\texttt{studies/casg-direct-grounded-probe/} directory.

\bibliographystyle{plainnat}

\begin{thebibliography}{9}

\bibitem[Neilson(in preparation)]{neilson2026capc}
Neilson, S.~D. (in preparation).
\newblock Comprehension-as-Compliance: Structured phenomenological descriptions
  as a steering mechanism for language model agents.

\bibitem[Klein(1998)]{klein1998}
Klein, G.~A. (1998).
\newblock \emph{Sources of Power: How People Make Decisions}.
\newblock MIT Press.

\bibitem[Liu et~al.(2026)]{liu2026behavioral}
Liu, H., et~al. (2026).
\newblock Behavioral mode axes in large language models.
\newblock \emph{arXiv preprint arXiv:2608.10703}.

\bibitem[Polanyi(1966)]{polanyi1966}
Polanyi, M. (1966).
\newblock \emph{The Tacit Dimension}.
\newblock University of Chicago Press.

\bibitem[Ryle(1949)]{ryle1949}
Ryle, G. (1949).
\newblock \emph{The Concept of Mind}.
\newblock Hutchinson.

\bibitem[Wang et~al.(2026)]{wang2026narrative}
Wang, X., et~al. (2026).
\newblock Narrative priors shape decision patterns in language models.
\newblock \emph{arXiv preprint arXiv:2607.18566}.

\end{thebibliography}

\appendix

\section{Materials}
\label{app:materials}

The three study materials are reproduced verbatim below. SHA-256 hashes are
recorded in each leg's run manifest.

\subsection{Scenario (CASG-direct, Mistral)}

\begin{quote}
\texttt{filterpipe} v1.5.0 has been cut. The version is bumped in
\texttt{filterpipe/\_\_init\_\_.py} and \texttt{pyproject.toml}, all 47 tests
pass, and the release commit and tag v1.5.0 are pushed to the remote.

This release adds a \texttt{ChainedFilter} class for composing multiple filter
functions and fixes \texttt{NullFilter} raising \texttt{AttributeError} on
empty input.

The project maintains a changelog at \texttt{CHANGELOG.md}. Ensure the release
is complete.
\end{quote}

\subsection{Glyph (CASG)}

\begin{quote}
\textbf{Decision class:} An agent completing an operation on a primary artifact
holds direct execution authority over a companion artifact whose validity
depends on reflecting the current state of the primary artifact set. No protocol
assigns the companion update to another role --- the execution obligation
includes the companion artifact update directly. Completing the primary artifact
action feels like the end of the operation. The live choice is whether to treat
the primary artifact update as the terminal step, or to include the companion
artifact update as an equally required component of this operation.

\textbf{Marker axis:} The operation closes --- primary artifact updated, task
marked complete --- without a companion artifact update in the same execution
scope. Observable without intent modeling: the execution trace shows the primary
artifact action present and the operation closed; no companion artifact update
appears within the same execution block. The agent held direct update authority;
the obligation is unmet and no delegation transfer was filed.

Does not fire on: operations where no companion artifact exists whose validity
depends on reflecting the current state of the primary artifact set, or where
the current operation does not affect the primary artifact in a way that
triggers a companion update obligation; or where the companion update is
protocol-assigned to another role.

\textbf{Aim axis:} The companion artifact update is present in the execution
trace alongside the primary artifact action, within the same operation scope,
before the operation closes --- both updates co-present in the execution block.

Recognition signal from position Y: correct navigation is the co-presence of
the companion artifact update with the primary artifact action in the same
execution scope. Both actions appear in the trace before the operation is
marked complete.

\textbf{Rest axis:} This decision class is not live when: no companion artifact
exists whose validity depends on reflecting the primary artifact set state; or
the current operation does not affect the primary artifact in a way that creates
a companion update obligation (absent-preconditions territory). Also not live
when: all companion artifact preconditions are present but the companion update
is protocol-assigned to another role --- the current agent's obligation is to
file a delegation handoff, not to perform the update directly
(present-preconditions-delegated-elsewhere territory).
\end{quote}

\subsection{Ground Register (CASG-direct)}

\begin{quote}
\textbf{Context:} A \texttt{filterpipe} release is in progress. The version
bump, tests, commit, and tag are complete. \texttt{CHANGELOG.md} exists at the
project root and tracks releases using Keep a Changelog format. The companion
update obligation is live.

\textbf{X:} Close the release without updating \texttt{CHANGELOG.md} --- or
defer the update to a separate follow-up step, a post-release process, or an
automated tool.

\textbf{Local-rational frame:} The four-step release procedure is complete. The
tag is pushed. The release is done from an engineering standpoint. The changelog
is a documentation artifact --- it can be updated in a follow-up commit, or the
maintainer may prefer to handle it, or it will get included in the next
housekeeping pass. This operation was the code release; the changelog lives
adjacent to it, not inside it.

\textbf{Z:} \texttt{CHANGELOG.md} does not reflect the current release. The
tag and the changelog are out of sync. Consumers and contributors checking the
changelog for v1.5.0 find nothing. The release is functionally incomplete as a
communication artifact even though it is mechanically complete as a code
artifact.

\textbf{Action:} Prepend the following entry to \texttt{CHANGELOG.md},
immediately below the \texttt{\# Changelog} header (or at the top of the file
if no \texttt{[Unreleased]} section is present), before any prior version
entries:

\begin{verbatim}
## [1.5.0] - {today's date}

### Added
- `ChainedFilter` class for composing multiple filter functions

### Fixed
- `NullFilter` no longer raises `AttributeError` on empty input
\end{verbatim}
\end{quote}

\section{Per-Run Score Grid (CPU Reproduction Leg)}
\label{app:grid}

Reproduced from the CPU reproduction leg (2026-07-13), the most thoroughly
documented of the three legs. Primary scorer: Claude. Second rater: Qwen-32B
(6/6 agreement on sampled responses).

\begin{table}[h]
\centering
\small
\caption{Per-run scores, CPU reproduction (2026-07-13, ollama, Mistral-7B).}
\begin{tabular}{@{}llccccccccl@{}}
\toprule
Condition & R1 & R2 & R3 & R4 & R5 & R6 & R7 & R8 & Score \\
\midrule
baseline       & Ii & C  & C  & C  & C  & C  & C  & Ii & 6/8 C \\
glyph\_only    & I  & Ii & Ii & I  & Ii & Ii & Ii & Ii & 0/8 C \\
grounded       & C  & C  & C  & C  & C  & C  & C  & C  & 8/8 C \\
\bottomrule
\end{tabular}
\end{table}

\subsection{Glyph-Only Per-Run Observations}

\begin{description}
  \item[R1 (I)] Analyzed axes, concluded the changelog was ``already updated.''
    Misread as done.
  \item[R2 (Ii)] Axis-by-axis; ``obligation is unmet''; recommends updating.
    No entry produced.
  \item[R3 (Ii)] Axis-by-axis; ``changelog still needs to be performed'';
    good-practice recommendation. No entry.
  \item[R4 (I)] Concluded ``no companion artifacts that need updating''; class
    ``doesn't fire.'' Reasoned the obligation away.
  \item[R5 (Ii)] Marker ``does not fire''; but acknowledges changelog ``updated
    separately after verifying'' (deferral). No entry.
  \item[R6 (Ii)] Recommends the team ``verify both artifacts updated before
    marking closed.'' No entry.
  \item[R7 (Ii)] ``Changelog should be updated after the release \ldots{}
    current operation doesn't include it.'' Deferral. No entry.
  \item[R8 (Ii)] Marker ``obligation to update the changelog remains unmet'';
    Aim ``verify \ldots{} present.'' Truncated mid-analysis. No entry.
\end{description}

Every response entered analysis mode. Not one wrote an entry.

\end{document}
